REFERENCE

[1] Shi, S., Gao, F., Meng, X., Xu, X., & Zhu, J. (2019). Improving facial attractiveness prediction via co-attention learning. In 2019 IEEE International Conference on Acoustics, Speech and Signal Processing (ICASSP) (pp. 4045-4049). IEEE.

[2] Lugaresi, C., Tang, J., Nash, H., McClanahan, C., Uboweja, E., Hays, M., ... & Grundmann, M. (2019). MediaPipe: A framework for building perception pipelines. arXiv preprint arXiv:1906.08172.

[3] Liu, S., Yang, J., Huang, C., & Yang, M. H. (2016). Multi-objective convolutional learning for face labeling. In Proceedings of the IEEE Conference on Computer Vision and Pattern Recognition (pp. 3451-3459).

[4] Guo, J., Zhu, X., Yang, Y., Yang, F., Lei, Z., & Li, S. Z. (2020). Towards fast, accurate and stable 3D dense face alignment. In Proceedings of the European Conference on Computer Vision (ECCV) (pp. 152-168).

[5] Zhang, S., Zhu, X., Lei, Z., Shi, H., Wang, X., & Li, S. Z. (2017). FaceBoxes: A CPU real-time face detector with high accuracy. In 2017 IEEE International Joint Conference on Biometrics (IJCB) (pp. 1-9). IEEE.

[6] Farkas, L. G. (1994). Anthropometry of the Head and Face (2nd ed.). Raven Press, New York.

[7] Legan, H. L., & Burstone, C. J. (1980). Soft tissue cephalometric analysis for orthognathic surgery. Journal of Oral Surgery, 38(10), 744-751.

[8] Arnett, G. W., & Bergman, R. T. (1993). Facial keys to orthodontic diagnosis and treatment planning. American Journal of Orthodontics and Dentofacial Orthopedics, 103(4), 299-312.

[9] McNamara, J. A. Jr. (1984). A method of cephalometric evaluation. American Journal of Orthodontics, 86(6), 449-469.

[10] Ricketts, R. M. (1981). Perspectives in the clinical application of cephalometrics. Angle Orthodontist, 51(2), 115-150.

[11] Steiner, C. C. (1953). Cephalometrics for you and me. American Journal of Orthodontics, 39(10), 729-755.

[12] Holdaway, R. A. (1983). A soft-tissue cephalometric analysis and its use in orthodontic treatment planning. American Journal of Orthodontics, 84(1), 1-28.

[13] Burstone, C. J. (1967). Lip posture and its significance in treatment planning. American Journal of Orthodontics, 53(4), 262-284.

[14] Krizhevsky, A., Sutskever, I., & Hinton, G. E. (2012). ImageNet classification with deep convolutional neural networks. In Advances in Neural Information Processing Systems (pp. 1097-1105).

[15] He, K., Zhang, X., Ren, S., & Sun, J. (2016). Deep residual learning for image recognition. In Proceedings of the IEEE Conference on Computer Vision and Pattern Recognition (pp. 770-778).

[16] Sandler, M., Howard, A., Zhu, M., Zhmoginov, A., & Chen, L. C. (2018). MobileNetV2: Inverted residuals and linear bottlenecks. In Proceedings of the IEEE Conference on Computer Vision and Pattern Recognition (pp. 4510-4520).

[17] Paszke, A., Gross, S., Massa, F., Lerer, A., Bradbury, J., Chanan, G., ... & Chintala, S. (2019). PyTorch: An imperative style, high-performance deep learning library. In Advances in Neural Information Processing Systems (pp. 8024-8035).

[18] Bradski, G. (2000). The OpenCV Library. Dr. Dobb's Journal of Software Tools.

[19] Face++. (2026). Face Detect API Documentation: Beauty Attribute. Retrieved June 23, 2026, from https://console.faceplusplus.com/documents/5679127

[20] Fotor. (2026). Golden Ratio Face Calculator: Upload Photo Face Score Online. Retrieved June 23, 2026, from https://www.fotor.com/features/golden-ratio-face-calculator/

[21] Face Technology Company. (2025). Beauty Scanner - Face Analyzer. Retrieved June 23, 2026, from https://play.google.com/store/apps/details?id=com.golden.ratio.face

[22] FaceRate AI. (2026). AI Face Rate: Free Online Attractiveness Test. Retrieved June 23, 2026, from https://aifacerate.com/

[23] Xu, J., Jin, L., Liang, L., Feng, Z., Xie, D., & Mao, H. (2017). Facial attractiveness prediction using psychologically inspired convolutional neural network (PI-CNN). In 2017 IEEE International Conference on Acoustics, Speech and Signal Processing (ICASSP) (pp. 1657-1661). IEEE.

[24] Gray, D., Yu, K., Xu, W., & Gong, Y. (2010). Predicting facial beauty without landmarks. In European Conference on Computer Vision (pp. 434-447). Springer, Berlin, Heidelberg.

[25] Eisenthal, Y., Dror, G., & Ruppin, E. (2006). Facial attractiveness: Beauty and the machine. Neural Computation, 18(1), 119-142.

[26] Kagian, A., Dror, G., Leyvand, T., Meilijson, I., Cohen-Or, D., & Ruppin, E. (2008). A machine learning predictor of facial attractiveness revealing human-like psychophysical biases. Vision Research, 48(2), 235-243.

[27] Schmid, K., Marx, D., & Samal, A. (2008). Computation of a face attractiveness index based on neoclassical canons, symmetry, and golden ratios. Pattern Recognition, 41(8), 2710-2717.

[28] Laurentini, A., & Bottino, A. (2014). Computer analysis of face beauty: A survey. Computer Vision and Image Understanding, 125, 184-199.

[29] Xie, D., Liang, L., Jin, L., Xu, J., & Li, M. (2015). SCUT-FBP: A benchmark dataset for facial beauty perception. In 2015 IEEE International Conference on Systems, Man, and Cybernetics (pp. 1821-1826). IEEE.

[30] Liang, L., Lin, L., Jin, L., Xie, D., & Li, M. (2018). SCUT-FBP5500: A diverse benchmark dataset for multi-paradigm facial beauty prediction. In 2018 24th International Conference on Pattern Recognition (ICPR) (pp. 1598-1603). IEEE.
